
%%%%%%%%%%%%%%%%%%%%%%%%%%%%%%%%%%%%%%%%%%%%%%%%%%%%%%%%%%%%%%%%%%%%%%%%%%%%%%%%
\section{Introduction}
\label{sec:introduction}
%%%%%%%%%%%%%%%%%%%%%%%%%%%%%%%%%%%%%%%%%%%%%%%%%%%%%%%%%%%%%%%%%%%%%%%%%%%%%%%%

\subsection{Background and Motivation}
\label{subsec:intro_background}

The rapid digital transformation of modern industry has fundamentally changed the role of wireless communication systems. Emerging paradigms such as Industry~4.0, the Industrial Internet of Things (\gls{iiot}), cyber-physical systems (\gls{cps}), intelligent manufacturing, autonomous robotics, digital twins, and edge computing increasingly rely on seamless connectivity among heterogeneous devices, sensors, actuators, and control systems. Unlike traditional communication networks that primarily support human-centric services, modern industrial environments require continuous machine-to-machine communication capable of supporting real-time monitoring, distributed control, predictive maintenance, collaborative robotics, and autonomous decision-making. Consequently, wireless communication has evolved from a convenient alternative to wired connectivity into a fundamental enabling technology for digital transformation across manufacturing, transportation, healthcare, logistics, agriculture, and other industrial sectors \cite{xuInternetThingsIndustries2014,sisinniIndustrialInternetThings2018,boyesIndustrialInternetThings2018,wollschlaegerFutureIndustrialCommunication2017a}.

The communication requirements of these emerging applications differ significantly from those of conventional mobile broadband services. Industrial communication systems demand deterministic medium access, ultra-high reliability, low end-to-end latency, synchronized operation among distributed devices, massive machine connectivity, high energy efficiency, and predictable communication performance even under harsh electromagnetic conditions. Simultaneously, wireless infrastructures must support large numbers of heterogeneous devices while maintaining scalability, deployment flexibility, and long operational lifetimes with minimal maintenance. These often conflicting requirements have become one of the principal driving forces behind the evolution of next-generation wireless communication technologies \cite{liIndustrialInternetSurvey2017,sisinniIndustrialInternetThings2018,xuInternetThingsIndustries2014}.

Although existing wireless technologies such as Wi-Fi, Bluetooth, IEEE~802.15.4, WirelessHART, and conventional cellular systems have enabled numerous industrial applications, each has been designed to optimize only a subset of these communication requirements. High-throughput technologies frequently struggle to provide deterministic communication and predictable latency, whereas low-power wireless networks often sacrifice scalability, mobility, or spectral efficiency. Similarly, traditional cellular systems provide extensive coverage and mobility support but generally rely on licensed spectrum and infrastructure-centric deployments that may not be suitable for all industrial scenarios. These limitations have motivated the development of new radio access technologies capable of combining the flexibility of license-exempt operation with the reliability, scalability, and quality-of-service required by next-generation industrial communication systems \cite{wollschlaegerFutureIndustrialCommunication2017a,boccardiFiveDisruptiveTechnology2014,agiwalNextGeneration5G2016}.

%%%%%%%%%%%%%%%%%%%%%%%%%%%%%%%%%%%%%%%%%%%%%%%%%%%%%%%%%%%%%%%%%%%%%%%%%%%%%%%%
\subsection{Evolution toward IMT-2020 Wireless Communications}
\label{subsec:intro_imt}
%%%%%%%%%%%%%%%%%%%%%%%%%%%%%%%%%%%%%%%%%%%%%%%%%%%%%%%%%%%%%%%%%%%%%%%%%%%%%%%%

The evolution of wireless communication has been characterized by continuous improvements in capacity, spectral efficiency, mobility support, and service diversity. While second- and third-generation cellular systems primarily focused on voice and mobile data services, fourth-generation technologies introduced all-IP architectures capable of supporting broadband multimedia applications. However, the rapid emergence of industrial automation, cyber-physical systems, and large-scale machine communications exposed new requirements that could not be fully addressed through incremental improvements of existing cellular technologies. Future wireless networks were expected to simultaneously support enhanced broadband services, latency-critical industrial applications, and massive deployments of connected devices, creating the need for a fundamentally new communication framework \cite{andrewsWhatWill5G2014a,boccardiFiveDisruptiveTechnology2014,osseiranScenarios5GMobile2014,agiwalNextGeneration5G2016,parkvallNRNew5G2017a,dahlman2013lteadvanced}.

To address these challenges, the International Telecommunication Union Radiocommunication Sector (ITU-R) established the International Mobile Telecommunications-2020 (\gls{imt2020}) framework, defining a comprehensive vision for future wireless communication systems. Rather than specifying a single radio technology, IMT-2020 establishes performance objectives, minimum technical requirements, evaluation methodologies, and interoperability criteria that candidate radio interface technologies must satisfy. The framework defines three complementary service categories, namely enhanced Mobile Broadband (\gls{embb}), Ultra-Reliable Low-Latency Communications (\gls{urllc}), and massive Machine-Type Communications (\gls{mmtc}), which together extend wireless communication far beyond conventional mobile broadband toward industrial automation, mission-critical networking, intelligent transportation, and massive IoT deployments \cite{itur_m2083,itur_m2410,itur_m2412,itur_m2150}.

An important characteristic of \gls{imt2020} is its technology-neutral philosophy. Unlike previous generations, which were largely associated with a single family of cellular technologies, \gls{imt2020} establishes a common performance, evaluation, and interoperability framework that enables multiple \glspl{rit} to satisfy the same set of technical requirements rather than prescribing a single air interface \cite{itur_m2083,itur_m2410,itur_m2412,itur_m2150}. Although 3GPP New Radio (NR) represents the dominant cellular implementation within the \gls{imt2020} ecosystem, the framework also accommodates alternative radio technologies designed for deployment scenarios where conventional cellular solutions may not provide the most appropriate balance between performance, deployment flexibility, and spectrum utilization.

Within this context, \gls{etsi} initiated the development of \gls{dectnr} as a next-generation wireless technology operating in globally available license-exempt spectrum. Building upon the legacy of the original \gls{dect} standard while introducing a completely new radio interface, \gls{dectnr} was designed to support reliable, scalable, and deterministic wireless communication for industrial automation, massive machine-type communications, professional audio, and other mission-critical applications. Following a comprehensive evaluation against the technical performance requirements defined by \gls{imt2020}, \gls{dectnr} was officially approved by the \gls{itur} as a \gls{rit}, becoming the first \gls{etsi}-developed non-cellular radio interface included in the \gls{imt2020} family \cite{etsi_tr103810,etsi_tr103943,maliatsosEvaluationIMT2020Candidate2022a,perez-guiraoDECTNRUnveiling2022,kovalchukovDECT2020NewRadio2022}. This achievement significantly broadened the technological diversity of \gls{imt2020} by introducing a wireless communication platform specifically optimized for industrial, enterprise, and other license-exempt deployment scenarios, thereby complementing existing cellular solutions rather than replacing them \cite{smidaETSIDECT2020Component2023,haqueDECT2020NRLink2025}.
%%%%%%%%%%%%%%%%%%%%%%%%%%%%%%%%%%%%%%%%%%%%%%%%%%%%%%%%%%%%%%%%%%%%%%%%%%%%%%%%
\subsection{DECT-2020 NR: A New IMT-2020 Radio Interface}
\label{subsec:intro_dect}
%%%%%%%%%%%%%%%%%%%%%%%%%%%%%%%%%%%%%%%%%%%%%%%%%%%%%%%%%%%%%%%%%%%%%%%%%%%%%%%%

Among the radio interface technologies accepted within the IMT-2020 framework, DECT-2020 New Radio (DECT-2020 NR) represents a unique milestone in the evolution of wireless communications. Unlike conventional cellular technologies, DECT-2020 NR was specifically developed by the European Telecommunications Standards Institute (ETSI) to provide a flexible, scalable, and reliable wireless communication system operating in globally available license-exempt spectrum. Building upon the successful legacy of Digital Enhanced Cordless Telecommunications (DECT), the new radio interface extends the technology far beyond cordless telephony by introducing a modern packet-oriented architecture capable of supporting industrial automation, massive machine-type communications, ultra-reliable low-latency services, and mission-critical wireless networking. Its acceptance by ITU-R as an official IMT-2020 radio interface established DECT-2020 NR as the first non-cellular technology recognized within the IMT-2020 family, significantly broadening the technological diversity of future wireless communication systems \cite{itur_m2150,etsi_tr103810,maliatsosEvaluationIMT2020Candidate2022a,perez-guiraoDECTNRUnveiling2022}.

The standardization of DECT-2020 NR was carried out through a comprehensive family of ETSI Technical Specifications defining the overall system architecture, protocol stack, physical layer, medium access control procedures, convergence layer, and implementation framework. Release~1 established the initial protocol architecture and radio interface, while Release~2 introduced numerous enhancements aimed at improving system flexibility, scalability, interoperability, and industrial applicability. Together with the accompanying ETSI Technical Reports, these specifications provide a complete framework covering protocol operation, spectrum usage, regulatory aspects, coexistence mechanisms, and evaluation procedures required for IMT-2020 compliance \cite{etsi_ts103636_1,etsi_ts103636_2,etsi_ts103636_3,etsi_ts103636_4,etsi_ts103636_5,etsi_tr103943,etsi_tr103810}.

Unlike conventional cellular deployments that depend on licensed spectrum and centralized network infrastructures, DECT-2020 NR emphasizes autonomous network operation, decentralized resource management, efficient spectrum utilization, and flexible deployment in industrial and enterprise environments. These characteristics make the technology particularly attractive for private wireless networks, industrial Internet of Things applications, cyber-physical production systems, warehouse automation, autonomous mobile robots, smart logistics, and other mission-critical industrial scenarios where deterministic communication, low deployment cost, and spectrum independence are essential. Consequently, DECT-2020 NR has attracted increasing attention from both academia and industry as a complementary wireless technology capable of addressing communication scenarios that remain challenging for conventional cellular systems \cite{kovalchukovDECT2020NewRadio2022,perez-guiraoDECTNRUnveiling2022,smidaETSIDECT2020Component2023}.
%%%%%%%%%%%%%%%%%%%%%%%%%%%%%%%%%%%%%%%%%%%%%%%%%%%%%%%%%%%%%%%%%%%%%%%%%%%%%%%%
\subsection{Evolution of the DECT-2020 NR Research Landscape}
\label{subsec:intro_landscape}
%%%%%%%%%%%%%%%%%%%%%%%%%%%%%%%%%%%%%%%%%%%%%%%%%%%%%%%%%%%%%%%%%%%%%%%%%%%%%%%%

Following the publication of the ETSI specifications, research on DECT-2020 NR rapidly evolved from standardization activities toward practical implementation and experimental validation. Early studies concentrated on validating the protocol architecture and investigating the fundamental characteristics of the physical and medium access control layers, including synchronization procedures, waveform generation, channel estimation, modulation and coding schemes, random access mechanisms, coexistence strategies, and radio resource management. These investigations confirmed the feasibility of the standardized design while establishing the theoretical and experimental foundations required for practical deployment \cite{smidaETSIDECT2020Component2023,haqueDECT2020NRLink2025,wassmannPerformanceDECT2020NR2024,glazkovImprovingNearURLLCServices2025}.

As implementation platforms became available, the focus of the research community progressively shifted toward hardware realization and system-level experimentation. Software-defined radio platforms, FPGA-based implementations, commercial Nordic Semiconductor development kits, embedded communication modules, and industrial testbeds enabled comprehensive experimental investigations under realistic deployment conditions. These studies evaluated protocol behavior, link performance, latency, throughput, synchronization accuracy, coexistence performance, and energy consumption, thereby demonstrating the practical viability of DECT-2020 NR for industrial wireless communication systems \cite{fagerholmFPGAbasedDECT2020New2024,bandaraExperimentalValidationAnalysis2025,jahangirAnalysisPerformanceDECT2024,grafParetoOptimalityApproach2026a,grafWhatExpectWhen2026}.

More recently, the research landscape has expanded considerably beyond the lower communication layers. Current investigations encompass scheduling and radio resource allocation, routing and multi-hop communication, mesh networking, positioning and localization, authentication and security, coexistence with heterogeneous wireless technologies, software-defined networking, energy-efficient communication, multi-RAT integration, industrial automation, professional audio applications, and emerging non-terrestrial communication scenarios. This diversification illustrates the rapid transition of DECT-2020 NR from a newly standardized radio interface into a comprehensive research field spanning wireless communications, embedded systems, industrial networking, and cyber-physical systems \cite{soniPerformanceMultihopAssisted2021,soniLinkLevelAssessmentNOMA2023,saikkoPositioningDECT2020NR2024,saikkoPositioningTrackingDECT20202025,fuhrlingRobustRoutingProtocol2025b,nihtilaEnergySavingRouter2022,dangUniRATUnifiedDECT20202026,fedrecheskiMariResponsiveWireless2026,piresDECT2020NREnergy2026b}.

%%%%%%%%%%%%%%%%%%%%%%%%%%%%%%%%%%%%%%%%%%%%%%%%%%%%%%%%%%%%%%%%%%%%%%%%%%%%%%%%
\subsection{Research Gap }
\label{subsec:intro_gap}
%%%%%%%%%%%%%%%%%%%%%%%%%%%%%%%%%%%%%%%%%%%%%%%%%%%%%%%%%%%%%%%%%%%%%%%%%%%%%%%%

The rapid evolution of \gls{dectnr} has resulted in a substantial increase in both the volume and diversity of the available scientific literature. As demonstrated by the systematic review presented in this paper, research has progressively evolved from standardization activities toward protocol optimization, hardware implementation, experimental validation, and industrial deployment. Consequently, the literature now encompasses a broad spectrum of research topics, including system architecture, \gls{phy} techniques, \gls{mac} protocols, synchronization, scheduling and resource allocation, routing and mesh networking, localization and positioning, security, coexistence mechanisms, embedded implementations, performance evaluation, and application-specific deployments. This diversification demonstrates the growing technological maturity of \gls{dectnr} while simultaneously increasing the complexity of understanding the overall research landscape \cite{etsi_ts103636_1,etsi_ts103636_2,etsi_ts103636_3,etsi_ts103636_4,etsi_ts103636_5,smidaETSIDECT2020Component2023,haqueDECT2020NRLink2025,wassmannPerformanceDECT2020NR2024,bandaraExperimentalValidationAnalysis2025,jahangirAnalysisPerformanceDECT2024,grafWhatExpectWhen2026,soniPerformanceMultihopAssisted2021,saikkoPositioningDECT2020NR2024,dangUniRATUnifiedDECT20202026}.

The findings of our systematic literature review further indicate that the available knowledge is distributed across multiple publication types and research communities. Fundamental aspects of the technology are documented in ETSI Technical Specifications, Technical Reports, and ITU-R recommendations, whereas implementation methodologies, protocol optimizations, hardware platforms, experimental evaluations, and industrial applications are primarily reported in journal articles, conference proceedings, industrial demonstrations, and experimental studies. Although these publications collectively provide comprehensive coverage of individual research topics, they are largely organized according to specific technical objectives rather than the overall evolution of the technology. Consequently, researchers and practitioners must consult numerous heterogeneous sources to obtain a comprehensive understanding of \gls{dectnr} and its current state of development \cite{itur_m2083,itur_m2150,etsi_tr103810,etsi_tr103943,etsi_ts103636_1,etsi_ts103636_2,etsi_ts103636_3,etsi_ts103636_4,etsi_ts103636_5,smidaETSIDECT2020Component2023,perez-guiraoDECTNRUnveiling2022,haqueDECT2020NRLink2025,bandaraExperimentalValidationAnalysis2025,grafParetoOptimalityApproach2026a}.

Furthermore, the continuous evolution of the standard itself has accelerated the diversification of research. Following the publication of the initial Release~1 specifications, Release~2 introduced additional capabilities that stimulated new investigations in radio resource management, synchronization, localization, industrial automation, energy-efficient communication, multi-hop networking, and emerging Industry~4.0 applications. At the same time, the availability of commercial chipsets, software development kits, \gls{fpga} implementations, and industrial testbeds has shifted the research focus from theoretical protocol design toward comprehensive experimental validation under realistic deployment conditions. As a result, the literature has become increasingly fragmented across protocol layers, implementation platforms, optimization techniques, and application domains \cite{etsi_ts103636_1,etsi_ts103636_2,etsi_ts103636_3,fagerholmFPGAbasedDECT2020New2024,haqueDECT2020NRLink2025,bandaraExperimentalValidationAnalysis2025,saikkoPositioningTrackingDECT20202025,nihtilaEnergySavingRouter2022,fedrecheskiMariResponsiveWireless2026,piresDECT2020NREnergy2026b}.

Our analysis of the selected literature further reveals that most existing publications investigate individual aspects of \gls{dectnr} in isolation. While these studies have significantly advanced the technology, relatively few establish relationships among different protocol layers, compare alternative technical approaches across multiple research domains, or synthesize the overall evolution of the technology from standardization to industrial deployment. Consequently, common design principles, emerging research trends, remaining technical challenges, and future research opportunities remain dispersed throughout the literature rather than being presented within a unified analytical framework \cite{smidaETSIDECT2020Component2023,haqueDECT2020NRLink2025,soniPerformanceMultihopAssisted2021,saikkoPositioningDECT2020NR2024,dangUniRATUnifiedDECT20202026,bandaraExperimentalValidationAnalysis2025}.

Motivated by these observations, this paper presents a comprehensive systematic literature review of \gls{dectnr} following the \gls{prisma}~2020 methodology. Unlike conventional narrative surveys, our review employs a transparent and reproducible protocol for literature identification, screening, quality assessment, data extraction, and qualitative synthesis. The selected publications are systematically classified into a unified taxonomy covering the complete \gls{dectnr} technology stack, including standardization, system architecture, \gls{phy}, \gls{mac}, scheduling and resource allocation, routing and mesh networking, localization, security, implementation platforms, experimental validation, performance evaluation, and industrial applications. By integrating these diverse research directions into a coherent analytical framework, this survey provides a comprehensive overview of the current state of \gls{dectnr} research while identifying technological trends, remaining research challenges, and promising future research directions.

%%%%%%%%%%%%%%%%%%%%%%%%%%%%%%%%%%%%%%%%%%%%%%%%%%%%%%%%%%%%%%%%%%%%%%%%%%%%%%%%
\subsection{Contributions of This Survey}
\label{subsec:intro_contributions}
%%%%%%%%%%%%%%%%%%%%%%%%%%%%%%%%%%%%%%%%%%%%%%%%%%%%%%%%%%%%%%%%%%%%%%%%%%%%%%%%

The main contributions of this survey are summarized as follows:

\begin{enumerate}

\item We present the first comprehensive \gls{prisma}-based systematic literature review of DECT-2020 NR research, employing a transparent, reproducible, and evidence-based methodology for literature identification, screening, quality assessment, data extraction, and qualitative synthesis.

\item We provide a unified taxonomy that systematically classifies DECT-2020 NR research into coherent technical categories spanning the complete technology stack, including standardization, system architecture, physical layer, medium access control, scheduling and resource allocation, routing and mesh networking, localization, security, implementation platforms, experimental validation, performance evaluation, and industrial applications.

\item We critically synthesize the existing literature by integrating ETSI standards, ITU-R recommendations, journal articles, conference papers, and industrial implementation studies, thereby providing a comprehensive overview of the technological evolution of DECT-2020 NR from standardization to practical deployment.

\item We analyze the current state of research across the identified technical categories, highlighting the major technological developments, implementation approaches, experimental achievements, and emerging research trends that characterize the evolution of DECT-2020 NR.

\item We identify the remaining research challenges, unresolved technical issues, and promising future research directions, providing guidance for researchers, system developers, standardization bodies, and industrial practitioners working on next-generation industrial wireless communication systems.

\end{enumerate}
%%%%%%%%%%%%%%%%%%%%%%%%%%%%%%%%%%%%%%%%%%%%%%%%%%%%%%%%%%%%%%%%%%%%%%%%%%%%%%%%
\subsection{Paper Organization}
\label{subsec:intro_organization}
%%%%%%%%%%%%%%%%%%%%%%%%%%%%%%%%%%%%%%%%%%%%%%%%%%%%%%%%%%%%%%%%%%%%%%%%%%%%%%%%

The remainder of this paper is organized as follows.
Section~\ref{sec:methodology} presents the systematic review methodology, including the search strategy, study selection, quality assessment, data extraction, and qualitative synthesis procedures.
Section~\ref{sec:results} summarizes the literature search results, publication trends, and bibliometric analysis.
Section~\ref{sec:technical_synthesis} provides a comprehensive technical synthesis of the selected studies according to the proposed taxonomy.
Section~\ref{sec:future} discusses the remaining research challenges and outlines promising future research directions.
Finally, Section~\ref{sec:conclusion} concludes the paper.
